\section{Related work and positioning}
\label{sec:related-work}

\subsection{LPV modeling and vehicle control}
LPV models represent dynamics whose coefficients depend on measurable scheduling
variables, providing a systematic alternative to independently designed frozen
controllers. Established LPV references cover identification, embedding choices,
gain scheduling, and stability analysis
\cite{toth2010lpv,apkarian1995selfscheduled,hoffmann2015survey}. Vehicle-control
studies have used LPV formulations to expose velocity dependence and tire
nonlinearities to scheduled controllers, including path-tracking designs validated
in high-fidelity simulation \cite{corno2021lpv,gimondi2021lpv}. The bicycle and
tire models used here follow the standard vehicle-dynamics hierarchy from linear
cornering stiffness to nonlinear saturation \cite{pacejka2012tire}.

This project does not claim novelty for representing vehicle dynamics as LPV or
for interpolating gains over speed. Its LPV contribution is representational:
known within-client operating variation is assigned to scheduling, so that
persistent between-client differences can be learned without proliferating
unrelated frozen-speed models.

\subsection{Federated learning under heterogeneity}
Federated averaging established the shared-model paradigm in which clients retain
their raw data and communicate locally computed updates \cite{mcmahan2017fedavg}.
When client objectives differ, federated multi-task and clustered methods replace
a single compromise model with structured sharing
\cite{smith2017federated,ghosh2020ifca,sattler2021clustered}. These methods motivate
family-specific aggregation, but generally treat heterogeneity as a statistical
or task-level property rather than separating a measured scheduling coordinate
from a persistent physical phenotype.

\subsection{Federated system identification and the research gap}
FedSysID studies collaborative identification when clients observe similar but
non-identical linear dynamical systems and quantifies the tradeoff between added
samples and system heterogeneity \cite{wang2023fedsysid}. This directly supports
the premise that a fleet can improve system identification through collaboration,
while also showing why unstructured global aggregation can become harmful as
physical differences increase.

The gap addressed here lies at the intersection of these bodies of work. Existing
LPV vehicle-control studies primarily assume a prescribed plant or centralized
identification; clustered FL methods learn multiple client models but do not
encode continuous operating-point variation; and current federated system
identification centers on shared or nearby fixed-system models. The proposed
architecture instead seeks a small number of \emph{LPV model manifolds}: each
family shares a speed-dependent model, while family membership captures persistent
vehicle heterogeneity. The oracle experiments in
\cref{sec:phase0,sec:phase1,sec:phase2,sec:phase3} first test whether this
decomposition is useful before any claim is made about recovering it from
decentralized data.
